\section{Abordagem}
\label{sec:abordagem}

Esta é a seção central do artigo e normalmente ocupa de três a cinco páginas no
formato SBC. Descreva a solução proposta com detalhe suficiente para que outra
pessoa consiga reproduzi-la, organizando o conteúdo em subseções quando isso
facilitar a leitura.

\subsection{Visão Geral}
\label{sec:abordagem:visao}

Apresente primeiro a ideia em alto nível, com uma figura ou um diagrama que
mostre como as partes se encaixam. Figuras e tabelas bem escolhidas economizam
prosa e tornam o raciocínio mais fácil de seguir.

\subsection{Detalhamento}
\label{sec:abordagem:detalhe}

Em seguida, aprofunde cada componente, mantendo o padrão de conceito, fonte de
apoio e instância concreta. Equações numeradas ajudam a formalizar definições,
como na Equação~\ref{eq:exemplo}.

\begin{equation}
  \label{eq:exemplo}
  f(x) = \sum_{i=1}^{n} w_i \, x_i + b
\end{equation}
